\section{Proposal}\label{sec:proposal}

\begin{frame}{\secname}
	\framesubtitle{Macro-based Specialization}

	Employ the powerful metaprogramming capabilities of Rust to build a \alert{macro library} that simulates native specialization, with:

	\begin{tabular}{p{0.45\textwidth} p{0.55\textwidth}}
		\begin{itemize}
			\item Native-like look and feel % similar syntax to native specialization, easy to use and easy to migrate existing code
			\item \alert{Compile-time resolution} % all the monomorphized implementations are available and the method calls are resolved at compile time 
			\item Sound and stable % no unsafe code, no compiler modifications, sits on top of stable Rust
		\end{itemize}
		 &
		\begin{itemize}
			\item \alert{Monomorphization} % no vtables, no dynamic dispatch, because everything is monomorphized at compile time
			\item Powerful API % predicates support (all, any, not) and support for most common use cases
		\end{itemize}
	\end{tabular}

	\pause
	\smallskip

	To do this, we use two macros:
	\texttt{when} % which is an attribute macro applied to the specialized implementations to specify their applicability conditions
	and \texttt{spec}. % which is a function-like macro wraps the method calls to resolve them at compile time

	\pause
	\bigskip

	The approach can also be applied to other languages % with powerful metaprogramming capabilities, needed to implement the macros
	or even directly into the compiler frontend. % since the compiler already has all the necessary information during the compilation process and can simply replicate the macro's behavior natively
	% TODO: frontend or backend?

\end{frame}